\section{The loop with Chapter~6, closed}
\label{sec:deamloop}

Here is the circle we have been walking around.  To run the tree, we
need the carry --- Chapter~6's pair $(F,D)$, the forward and the
discount factor.  Chapter~6 fitted them from put--call parity.  But
parity, as derived in \cref{sec:fwdtwo}, is a statement about
\emph{European} portfolios, and the quotes on the line were American:
each mid shifted by its own premium, more on the put side than the
call side, which is exactly why the residuals bowed and the naive
fitted forward sat \MacDeamBowNaiveGapBp\ bp low.  (Basis points of
\emph{forward} in this section, not of volatility --- we are back on
Chapter~6's ground.)  So the subtraction needs the forward, and the
forward needs the subtraction.  Neither comes first.

The resolution is the standard one for a fixed point: start with the
biased answer and iterate, because each half-step improves the input
to the other.  Concretely, on the running node's board of
\MacDeamBowNPairs\ paired strikes:
\begin{enumerate}[leftmargin=1.8em,itemsep=2pt]
\item \emph{Pass one.}  Take the naive fitted forward
$\hat F=\MacDeamBowFNaive$ from \cref{fig:fwdline} as the starting
carry.  De-Americanize every quote on the board --- both legs of
every pair --- through \eqref{eq:deamroot}, and subtract each leg's
premium estimate \eqref{eq:deamhat} from its mid.  Refit the parity
line \emph{on the de-Americanized mids}.  The fitted forward moves.
\item \emph{Pass two.}  Rebuild the carry from the new forward and do
it again.  This time the fitted forward moves by only
\MacDeamBowPassDeltaBp\ bp --- the loop has converged for any purpose
in this book, and a third pass is not worth its arithmetic.
\end{enumerate}
One input is deliberately \emph{not} taken from the board: the rate.
Chapter~6 proved the line's slope is the poorly identified
coefficient, and this board makes the point personally --- its fitted
slope implies $r=\MacFwdLineRateNaivePct\%$, a negative rate in
mid-2026, which is noise wearing a suit.  Following
\cref{sec:fwdclamp}'s discipline, the rate is pinned by a disclosed
convention (a money-market $4.3\%$; appendix~\ref{app:deamprotocol})
and the \emph{level} --- the well-measured number --- carries the
information, through $q_d=r-\log(F/S)/t$.

Not every pair survives the loop, and the dropouts are old friends:
on \MacDeamBowNDropped\ of the \MacDeamBowNPairs\ pairs, the deep
in-the-money put sits on its intrinsic floor --- \cref{sec:deamsub}'s
no-root lane, live on a real board --- so no premium can be measured
there and the pair is left out of the refit.  The
\MacDeamBowNKept\ both-leg pairs that remain are the population of
\cref{fig:deambow}.

\begin{figure}[htbp]
\centering
\paperfig{\textwidth}{fig_deam_bow.pdf}
\caption{Chapter~6's debt, paid (SPY December 2026,
\MacDeamBowNKept\ of \MacDeamBowNPairs\ pairs; the rest sit on the
intrinsic floor).  (a)~The measured parity residuals against the
tree's premium difference, both detrended by their own line: the dots
sit on the curve to \$\MacDeamBowMatchRmsDollars\ rms --- the bow
\emph{is} the premium difference.  (b)~Residuals after subtracting
each leg's premium: rms \$\MacDeamBowRmsBeforeDollars\ down to
\$\MacDeamBowRmsAfterDollars, and the refitted forward (blue) lands
\MacDeamBowGapToResolvedBp\ bp from the snapshot's stored resolved
forward (dotted) --- against \MacDeamBowNaiveGapBp\ bp for the naive
one (dashed).}
\label{fig:deambow}
\end{figure}

Panel~(a) settles an account left open since \cref{fig:fwdline}(b).
The orange dots are the measured residuals of the parity line on the
kept pairs.  The blue curve is \emph{the model's prediction of them}:
the tree's premium difference, call leg minus put leg, at each pair's
own fitted volatilities.  It is detrended by its own least-squares
line, exactly as the measured residuals are, because a parity
regression absorbs anything affine --- only the arch can survive it.
(For the same reason the arch here is shallower than
\cref{fig:fwdline}(b)'s: the line is refitted on the kept pairs
alone, which reshuffles how much of the bow the affine part absorbs;
the shape, not the split, carries the information.)  The dots
sit on the curve, to an rms of \$\MacDeamBowMatchRmsDollars\ --- the
same size as the quote noise left over in panel~(b), which is to say:
to within what this board can measure at all, the bow is
\emph{entirely} the premium difference.  The shape now reads itself.
At low strikes the calls are deep in the money, but SPY's dividend
carry is mild against the rate, so the call premium stays small ---
Merton's logic with a small leak.  At high strikes the put premium
grows dollar-deep and drags $\mathsf{C}-\mathsf{P}$ \emph{down};
after the fitted line absorbs the overall tilt, what remains is
precisely the arch Chapter~6 could describe but not explain.

Panel~(b) completes the repair.  With each leg's premium subtracted,
the residuals collapse from \$\MacDeamBowRmsBeforeDollars\ rms to
\$\MacDeamBowRmsAfterDollars\ --- from a structured bow, dollars
deep, to unstructured cent-scale noise --- and the refitted forward
moves from $\MacDeamBowFNaive$ to $\MacDeamBowFDeam$.  That lands
\MacDeamBowGapToResolvedBp\ bp from $\MacDeamBowFResolved$, the
forward stored with the book's frozen snapshot --- the resolution the
reference implementation (Chapter~6's working software, whose
conventions the book's figures audit) computed with its own discrete
dividend schedule.  Read that residual gap for what it is: the
distance between two disclosed dividend conventions
(\cref{sec:deamdivs}), not an unexplained error --- and compare it
with the \MacDeamBowNaiveGapBp\ bp the naive forward was carrying.
Chapter~6 measured, on this very node, what such basis points cost if
ignored: \MacFwdSkewGapBp\ vol bp of put--call gap at the money for
every \MacFwdSkewBumpBp\ bp of forward error, so the naive forward
would imprint \MacFwdSkewNaiveGapBp\ vol bp of artificial smile
structure (\cref{sec:fwdskew}'s own number).  The loop is not
bookkeeping; it is where the smile's raw material gets clean.

One extension is beyond this chapter, and Chapter~6 already flagged
it: on a genuinely hard-to-borrow name, the borrow $b$ enters the
carry, the carry moves the premium, and the premium moves the fitted
forward that identifies the borrow --- \cref{sec:fwdborrow}'s closing
caveat.  The resolution is this section's loop with one more unknown
riding inside it, and it inherits \cref{sec:fwdborrow}'s
identifiability floors --- the smallest borrow a board's noise lets
it distinguish from zero --- unchanged: below the floor, the jointly
closed answer is still \emph{unidentified}, reported as such.  On ordinary
names --- the frozen board included --- the borrow leg is immaterial
and the two passes above are the whole story.
